\chapter{The sound areas}\label{ch:sound}

\textsc{normal sound} and \textsc{drum sound} are the two areas of the address map that
hold an instrument sound rather than a setting. They are reached the same way as any other
parameter, with \textsc{itr} to write and \textsc{drq} to request.

\begin{center}
\begin{tabular}{lll}
\toprule
\textsc{adr} & address & area \\
\midrule
\bytes{10 00 00} & \bytes{040000} & \textsc{normal sound}, individual \\
\bytes{18 00 00} & \bytes{060000} & \textsc{drum sound}, individual \\
\bottomrule
\end{tabular}
\end{center}

\begin{caution}{This is the sound being edited, not a stored one}
These addresses reach the instrument's working copy --- the sound currently loaded for
editing. They are not a way to read or write one of the 32 stored sounds. For those, use
the \textsc{sound} bulk-dump category, which carries the whole sound memory at
\bytes{20 00 00}.
\end{caution}

\section{How a normal sound is laid out}
The \textsc{normal sound} area is \textbf{713 bytes}. It is one \textsc{general data}
block, then four identical \textsc{tone data} blocks, then four identical
\textsc{modeling data} blocks --- the instrument's four tones and the four driver and
resonator models that go with them.

\begin{center}
\begin{tabular}{lrrl}
\toprule
block & size & repeats & parameter numbers \\
\midrule
\textsc{general data}  & 217 & 1 & \bytes{000}--\bytes{0D8} \\
\textsc{tone data}     &  81 & 4 & \bytes{0D9}, \bytes{12A}, \bytes{17B}, \bytes{1CC} \\
\textsc{modeling data} &  43 & 4 & \bytes{21D}, \bytes{248}, \bytes{273}, \bytes{29E} \\
\midrule
total & 713 & & \bytes{000}--\bytes{2C8} \\
\bottomrule
\end{tabular}
\end{center}

A parameter number is an offset from the start of the area, so the address of parameter
\textit{p} of tone \textit{n} (counting tones from 1) is

\begin{center}
\bytes{040000} $+$ \bytes{0D9} $+$ 81 $\times$ (\textit{n} $-$ 1) $+$ \textit{p}
\end{center}

expressed as \textsc{adr} in the usual three seven-bit bytes, and the same with \bytes{21D}
and 43 for a modeling block. The tables below give one instance of each repeating block;
the other three are identical at the offsets above.

\begin{note}{Two independent sources agree on the size}
The block sizes and repeat counts here come from Technics' published parameter tables. The
figure of 713 bytes for the whole area comes from the program, decoded separately, and the
guide never states it. The transcribed blocks tile \bytes{000} to \bytes{2C8} with no gap
and no overlap and come to exactly that total, which is what makes the layout above
trustworthy rather than merely plausible.
\end{note}

\section{What is transcribed here, and what is not}
The tables give each parameter's number, its bit field within that byte, its name and its
range. They were read from scans of the published tables, and the numbering is checked by
the tiling described above. \textbf{The names, bit fields and ranges are not checked by
anything}; a misreading of one would survive.

Three enumerations the guide prints are not reproduced: the list of about fifty control
functions that a controller's \textsc{function select} byte chooses from, the meaning of
the filter \textsc{value} bytes for each filter mode, and the meaning of the twenty
\textsc{value} bytes under each effect type. All three are in the guide.

%% GENERATED by notes/sysex-probes/sysex_sound_layout.py --tex
%% Source: sound_layout.json. Do not edit.

\section{Normal Sound parameters}
\subsection*{GENERAL DATA}
{\small
\begin{longtable}{l>{\raggedright\arraybackslash}p{24mm}>{\raggedright\arraybackslash}p{58mm}>{\raggedright\arraybackslash}p{36mm}}
\toprule
no. & bits & parameter & values \\
\midrule\endfirsthead
\toprule no. & bits & parameter & values \\
\midrule\endhead
\bytes{000--00F} & 6--0 & NAME (16 characters) & \bytes{20-7F} --- space to \} \\
\bytes{010} & 7--6 & ATTRIBUTE FLAG: PARAMETER MODE & \bytes{00-00} --- 0: NORMAL MODE, fixed \\
\bytes{011} & 7--0 & TONE ON/OFF & two bits per tone, 4th 3rd 2nd 1st \\
\bytes{012} & 7--0 & KEY ON MODE & per tone: 0 KEY ON, 1 KEY OFF, 2 LEGATO, 3 NON LEGATO \\
\bytes{013} & 7--0 & KEY SCALING & \bytes{00-42} --- 0 OFF, 1-22 octave, 64-66 all key \\
\bytes{014--016} &  & CONTROLLERS: PITCH BEND & on/off and phase, function select, depth \\
\bytes{017--019} &  & MODULATION WHEEL 1 SELECT1 & same three bytes \\
\bytes{01A--01C} &  & MODULATION WHEEL 1 SELECT2 &  \\
\bytes{01D--01F} &  & MODULATION WHEEL 2 SELECT1 &  \\
\bytes{020--022} &  & MODULATION WHEEL 2 SELECT2 &  \\
\bytes{023} & 6--0 & REALTIME CREATOR BOX SELECT & bit per box, BOX1-BOX6 \\
\bytes{024} & 6--0 & REALTIME CONTROLLER BOX SELECT & bit per box, BOX1-BOX6 \\
\bytes{025--027} &  & CONTROL PEDAL SELECT1 &  \\
\bytes{028--02A} &  & CONTROL PEDAL SELECT2 &  \\
\bytes{02B--02D} &  & AFTER TOUCH SELECT1 &  \\
\bytes{02E--030} &  & AFTER TOUCH SELECT2 &  \\
\bytes{031--054} &  & BOX SELECT: CONTROLLER BOX1 X through BOX6 Y & twelve controllers, three bytes each, X then Y \\
\bytes{055} & 3--0 & PITCH OCTAVE SHIFT & \bytes{06-0A} --- 6 is -2 oct, 8 is 0, A is +2 oct \\
\bytes{056} &  & (RESERVE) &  \\
\bytes{057--066} &  & PITCH LFO BOX1-BOX4 & four bytes per box: depth, speed, key sync and touch depth, wave and delay \\
\bytes{067--076} &  & LEVEL LFO BOX1-BOX4 & same four bytes per box \\
\bytes{077--086} &  & FILTER LFO BOX1-BOX4 & same four bytes per box \\
\bytes{087} & 3--0, 7--4 & MIXER OUTPUT SELECT: MAIN, SUB & \bytes{00-05} --- 0 no output, 1 MAIN, 2-4 SUB1-3, 5 EFFECT2 \\
\bytes{088} & 6--0 & EFFECT SEND VOLUME: EFFECT1 & \bytes{00-7F} \\
\bytes{089} & 6--0 & EFFECT SEND VOLUME: REVERB & \bytes{00-7F} \\
\bytes{08A} & 3--0, 7--4 & EFFECT OUTPUT SELECT: EFFECT1, EFFECT2 & \bytes{01-04} --- 1 MAIN, 2-4 SUB1-3 \\
\bytes{08B} & 0 & DSP EFFECT COMMON ALGORITHM & \bytes{00-01} --- 0 SERIAL, 1 PARALLEL \\
\bytes{08C} &  & EFFECT1 TYPE & see the guide's DSP EFFECT pages \\
\bytes{08D--0A0} &  & EFFECT1 PARAMETER VALUE1-VALUE20 & meaning depends on the type \\
\bytes{0A1} &  & EFFECT1 ANOTHER EFF SEND &  \\
\bytes{0A2} &  & EFFECT1 DYNAMIC CONTROL &  \\
\bytes{0A3} &  & EFFECT2 TYPE &  \\
\bytes{0A4--0B7} &  & EFFECT2 PARAMETER VALUE1-VALUE20 &  \\
\bytes{0B8} &  & EFFECT2 ANOTHER EFF SEND &  \\
\bytes{0B9} &  & EFFECT2 DYNAMIC CONTROL &  \\
\bytes{0BA} &  & REVERB TYPE &  \\
\bytes{0BB--0CE} &  & REVERB PARAMETER VALUE1-VALUE20 &  \\
\bytes{0CF} &  & REVERB DYNAMIC CONTROL & reverb has no ANOTHER EFF SEND \\
\bytes{0D0} & 7, 6, 3--0 & DIGITAL EFFECT MODE: ON/OFF, PAN, EFFECT TYPE & \bytes{00-0B} --- 0-1 CELESTE, 2-3 CHORUS, 4-5 ENSEMBLE, 6-7 TREMOLO, 8-9 DELAY, 10-11 SOLO EFFECT \\
\bytes{0D1--0D8} &  & DIGITAL EFFECT VALUE1-VALUE8 &  \\
\bottomrule
\end{longtable}
}

\subsection*{TONE DATA}
{\small
\begin{longtable}{l>{\raggedright\arraybackslash}p{24mm}>{\raggedright\arraybackslash}p{58mm}>{\raggedright\arraybackslash}p{36mm}}
\toprule
no. & bits & parameter & values \\
\midrule\endfirsthead
\toprule no. & bits & parameter & values \\
\midrule\endhead
\bytes{0D9} & 5--0 & ATTRIBUTE DELAY & \bytes{00-32} --- 0 delay off \\
\bytes{0DA} & 7--0 & ATTRIBUTE PANNING & \bytes{00-80} --- 0-127 fixed pan, 128 random pan \\
\bytes{0DB} & 6--0 & TONE SELECT: TTN SELECT & \bytes{00-7F} \\
\bytes{0DC} & 7--6, 5--4, 3--0 & TONE KIND, TONE LOCATION, BANK SELECT & kind 0 NORMAL 1 DRUM 2 ATTACK; location 0 internal ROM 1 internal RAM 2 external ROM \\
\bytes{0DD} & 7--0 & PITCH KEY SHIFT & \bytes{E8-18} --- -24 to +24 \\
\bytes{0DE} & 7--0 & PITCH DETUNE & \bytes{80-7F} --- -128 to +127 \\
\bytes{0DF} & 7--6, 5, 4, 2--0 & BOX CONTROL LFO-FM1: BOX SELECT, ON/OFF, PHASE, SCALE SELECT & scale 0 1/1 through 6 1/64, 7 FIX \\
\bytes{0E0--0E9} &  & PITCH ENVELOPE SHAPE: total depth, start pitch, attack time, peak level, decay1 time, sustain1 level, decay2 time, sustain2 level, release time, stop pitch & one byte each, in that order \\
\bytes{0EA} & 7--0 & PITCH ENV TOUCH FOLLOW ADR TIME & \bytes{CE-32} --- -50 to +50 \\
\bytes{0EB} & 7--0 & PITCH ENV TOUCH FOLLOW DEPTH & \bytes{CE-32} \\
\bytes{0EC} & 6--0 & PITCH ENV KEY FOLLOW CENTER KEY & \bytes{00-7F} \\
\bytes{0ED} & 7--0 & PITCH ENV KEY FOLLOW ATTACK SLOPE & \bytes{CE-32} \\
\bytes{0EE} & 7--0 & PITCH ENV KEY FOLLOW DECAY SLOPE & \bytes{CE-32} \\
\bytes{0EF} & 7--0 & PITCH ENV KEY FOLLOW RELEASE SLOPE & \bytes{CE-32} \\
\bytes{0F0} & 6--0 & LEVEL VOLUME & \bytes{00-7F} --- one step is 0.375 dB \\
\bytes{0F1} & 7--0 & LEVEL TOUCH DEPTH & \bytes{CE-32} \\
\bytes{0F2} & 7--5 & LEVEL TOUCH CURVE & \bytes{00-06} --- 0 is -3 curve through 6 is +3 curve \\
\bytes{0F3} & 6--0 & LEVEL KEY FOLLOW CENTER KEY & \bytes{00-7F} \\
\bytes{0F4} & 6--0 & LEVEL KEY FOLLOW LOW KEY LIMIT & \bytes{00-7F} \\
\bytes{0F5} & 6--0 & LEVEL KEY FOLLOW HIGH KEY LIMIT & \bytes{00-7F} \\
\bytes{0F6} & 7--0 & LEVEL KEY FOLLOW SLOPE & \bytes{CE-32} \\
\bytes{0F7--0FA} & 6--0 & LAYER KEY RANGE: low key select, low fade, high key select, high fade & \bytes{00-7F} \\
\bytes{0FB--0FE} & 6--0 & LAYER VELOCITY RANGE: low vel select, low fade, high vel select, high fade & \bytes{00-7F} \\
\bytes{0FF} & 7--6, 5, 4 & BOX CONTROL LFO-AM: BOX SELECT, ON/OFF, PHASE &  \\
\bytes{100--106} &  & LEVEL ENVELOPE SHAPE: attack time, peak level, decay1 time, sustain1 level, decay2 time, sustain2 level, release time &  \\
\bytes{107} & 7--0 & LEVEL ENV TOUCH FOLLOW ATTACK & \bytes{CE-32} \\
\bytes{108} & 7--0 & LEVEL ENV TOUCH FOLLOW DECAY & \bytes{CE-32} \\
\bytes{109--10B} & 6--0 & LEVEL ENV KEY FOLLOW: center key, low key limit, high key limit & \bytes{00-7F} \\
\bytes{10C--10E} & 7--0 & LEVEL ENV KEY FOLLOW: attack slope, decay slope, release slope & \bytes{CE-32} \\
\bytes{10F} & 2--0, 7--5 & FILTER MODE, TOUCH CURVE & \bytes{00-06} --- 0 THROUGH, 1 LPF12+EQ, 2 HPF12+EQ, 3 LPF24, 4 HPF24, 5 BPF \\
\bytes{110} & 5--0 & FILTER TOUCH DEPTH & \bytes{00-32} \\
\bytes{111} & 7--6, 5, 4 & FILTER LFO-FCM: BOX SELECT, ON/OFF, PHASE &  \\
\bytes{112--114} & 6--0 & FILTER KEY FOLLOW: center key, low key limit, high key limit & \bytes{00-7F} \\
\bytes{115} & 7--0 & FILTER KEY FOLLOW SLOPE & \bytes{CE-32} \\
\bytes{116--11F} &  & FILTER ENVELOPE SHAPE: cutoff adjust, start point, attack time, peak level, decay1 time, sustain1 level, decay2 time, sustain2 level, release time, stop point &  \\
\bytes{120} & 7--0 & FILTER ENV TOUCH FOLLOW ADR TIME & \bytes{CE-32} \\
\bytes{121} & 7--0 & FILTER ENV TOUCH FOLLOW DEPTH & \bytes{CE-32} \\
\bytes{122} & 6--0 & FILTER ENV KEY FOLLOW CENTER KEY & \bytes{00-7F} \\
\bytes{123--125} & 7--0 & FILTER ENV KEY FOLLOW: attack slope, decay slope, release slope & \bytes{CE-32} \\
\bytes{126--129} &  & FILTER PARAMETER VALUE1-VALUE4 & meaning depends on the filter mode \\
\bottomrule
\end{longtable}
}

\subsection*{MODELING DATA}
{\small
\begin{longtable}{l>{\raggedright\arraybackslash}p{24mm}>{\raggedright\arraybackslash}p{58mm}>{\raggedright\arraybackslash}p{36mm}}
\toprule
no. & bits & parameter & values \\
\midrule\endfirsthead
\toprule no. & bits & parameter & values \\
\midrule\endhead
\bytes{21D--21F} & 6--0 & DRIVER THRESHOLD POINT: between A and B, B and C, C and D & \bytes{00-7F} --- provided that A < B < C < D \\
\bytes{220} & 6--0 & DRIVER A TTN SELECT & \bytes{00-7F} \\
\bytes{221} & 7--6, 5--4, 3--0 & DRIVER A KIND, LOCATION, BANK SELECT & kind 0 NORMAL 1 DRUM 2 ATTACK \\
\bytes{222--223} &  & DRIVER B: TTN select; kind, location, bank select &  \\
\bytes{224--225} &  & DRIVER C: TTN select; kind, location, bank select &  \\
\bytes{226--227} &  & DRIVER D: TTN select; kind, location, bank select &  \\
\bytes{228} & 7--6, 5--0 & RESONATOR GENERAL: INTERACTION CONNECT, RESONATOR TYPE & \bytes{00-3F} --- type 0 ORIGINAL, 1-63 preset \\
\bytes{229} &  & (RESERVE) &  \\
\bytes{22A} & 7--0 & POSITION PARAMETER POSITION & \bytes{00-FA} --- 0.2 per cent per step, 0 to 50 per cent \\
\bytes{22B} & 7, 6--0 & FORMANT FLAG, DEPTH & \bytes{00-7F} --- 0 FIXED, 1 MOVE \\
\bytes{22C} &  & (RESERVE) &  \\
\bytes{22D} & 7--0 & POSITION TOUCH & \bytes{CE-32} --- -50 to +50 \\
\bytes{22E} & 5--0 & POSITION MOVEMENT WIDTH & \bytes{00-32} \\
\bytes{22F} & 5--0, 7 & POSITION MOVEMENT SPEED, SAMPLE/HOLD & \bytes{00-32} \\
\bytes{230} & 6--0 & INTERACTION GAIN & \bytes{00-7F} \\
\bytes{231} &  & (RESERVE) &  \\
\bytes{232} & 7, 6--0 & MAIN RESONATOR MODE, FITTING & \bytes{00-7F} --- 0 positive feedback, 1 negative feedback \\
\bytes{233} & 7, 6--0 & MAIN RESONATOR SCALE, MUTING & \bytes{00-7F} \\
\bytes{234} & 7--0 & MAIN TOUCH FOLLOW TOUCH FITTING & \bytes{CE-32} \\
\bytes{235} & 7--0 & MAIN TOUCH FOLLOW TOUCH MUTING & \bytes{CE-32} \\
\bytes{236--238} & 6--0 & MAIN KEY FOLLOW: center key, low key limit, high key limit & \bytes{00-7F} \\
\bytes{239} & 7--0 & MAIN KEY FOLLOW SLOPE & \bytes{CE-32} \\
\bytes{23A} & 7--0 & MAIN PITCH KEY SHIFT & \bytes{C4-3C} --- -60 to +60 \\
\bytes{23B} & 7--0 & MAIN PITCH DETUNE & \bytes{80-7F} --- -128 to +127 \\
\bytes{23C} & 7, 6--0 & SUB RESONATOR MODE, FITTING & \bytes{00-7F} \\
\bytes{23D} & 7, 6--0 & SUB RESONATOR SCALE, MUTING & \bytes{00-7F} \\
\bytes{23E} & 7--0 & SUB GAIN & \bytes{9C-64} \\
\bytes{23F} & 7--0 & SUB TOUCH FOLLOW TOUCH FITTING & \bytes{CE-32} \\
\bytes{240} & 7--0 & SUB TOUCH FOLLOW TOUCH MUTING & \bytes{CE-32} \\
\bytes{241} & 7--0 & SUB TOUCH FOLLOW TOUCH SUB GAIN & \bytes{CE-32} \\
\bytes{242--244} & 6--0 & SUB KEY FOLLOW: center key, low key limit, high key limit & \bytes{00-7F} \\
\bytes{245} & 7--0 & SUB KEY FOLLOW SLOPE & \bytes{CE-32} \\
\bytes{246} & 7--0 & SUB PITCH KEY SHIFT & \bytes{C4-3C} \\
\bytes{247} & 7--0 & SUB PITCH DETUNE & \bytes{80-7F} \\
\bottomrule
\end{longtable}
}

\section{Drum Sound parameters}
\subsection*{KIT DATA}
{\small
\begin{longtable}{l>{\raggedright\arraybackslash}p{24mm}>{\raggedright\arraybackslash}p{58mm}>{\raggedright\arraybackslash}p{36mm}}
\toprule
no. & bits & parameter & values \\
\midrule\endfirsthead
\toprule no. & bits & parameter & values \\
\midrule\endhead
\bytes{000--00F} & 6--0 & NAME (16 characters) & \bytes{20-7F} --- space to \} \\
\bytes{010} & 7--6 & ATTRIBUTE FLAG: PARAMETER MODE & \bytes{00-00} --- 2: DRUM MODE, fixed \\
\bytes{011--013} &  & CONTROLLERS: PITCH BEND & on/off and phase, function select, depth \\
\bytes{014--016} &  & MODULATION WHEEL 1 SELECT1 & on/off and phase, function select, depth \\
\bytes{017--019} &  & MODULATION WHEEL 1 SELECT2 & on/off and phase, function select, depth \\
\bytes{01A--01C} &  & MODULATION WHEEL 2 SELECT1 & on/off and phase, function select, depth \\
\bytes{01D--01F} &  & MODULATION WHEEL 2 SELECT2 & on/off and phase, function select, depth \\
\bytes{020} & 6--0 & REALTIME CREATOR BOX SELECT & bit per box, BOX1-BOX6 \\
\bytes{021} & 6--0 & REALTIME CONTROLLER BOX SELECT & bit per box, BOX1-BOX6 \\
\bytes{022--024} &  & CONTROL PEDAL SELECT1 & on/off and phase, function select, depth \\
\bytes{025--027} &  & CONTROL PEDAL SELECT2 & on/off and phase, function select, depth \\
\bytes{028--02A} &  & AFTER TOUCH SELECT1 & on/off and phase, function select, depth \\
\bytes{02B--02D} &  & AFTER TOUCH SELECT2 & on/off and phase, function select, depth \\
\bytes{02E--051} &  & BOX SELECT: CONTROLLER BOX1 X through BOX6 Y & twelve controllers, three bytes each, X then Y \\
\bytes{052} & 3--0, 7--4 & MIXER EFFECT OUTPUT SELECT: EFFECT1, EFFECT2 & \bytes{01-04} --- 1 MAIN, 2-4 SUB1-3 \\
\bytes{053} & 0 & DSP EFFECT COMMON ALGORITHM & \bytes{00-01} --- 0 SERIAL, 1 PARALLEL \\
\bytes{054} &  & EFFECT1 TYPE & see the effect tables \\
\bytes{055--068} &  & EFFECT1 PARAMETER VALUE1-VALUE20 & meaning depends on the type \\
\bytes{069} &  & EFFECT1 ANOTHER EFF SEND &  \\
\bytes{06A} &  & EFFECT1 DYNAMIC CONTROL &  \\
\bytes{06B} &  & EFFECT2 TYPE &  \\
\bytes{06C--07F} &  & EFFECT2 PARAMETER VALUE1-VALUE20 &  \\
\bytes{080} &  & EFFECT2 ANOTHER EFF SEND &  \\
\bytes{081} &  & EFFECT2 DYNAMIC CONTROL &  \\
\bytes{082} &  & REVERB TYPE &  \\
\bytes{083--096} &  & REVERB PARAMETER VALUE1-VALUE20 &  \\
\bytes{097} &  & REVERB DYNAMIC CONTROL & reverb has no ANOTHER EFF SEND \\
\bottomrule
\end{longtable}
}

\subsection*{SOUND SELECT}
{\small
\begin{longtable}{l>{\raggedright\arraybackslash}p{24mm}>{\raggedright\arraybackslash}p{58mm}>{\raggedright\arraybackslash}p{36mm}}
\toprule
no. & bits & parameter & values \\
\midrule\endfirsthead
\toprule no. & bits & parameter & values \\
\midrule\endhead
\bytes{098--197} &  & SOUND SELECT, NOTE:00 to NOTE:7F & 128 notes, two bytes each \\
\bottomrule
\end{longtable}
}

\subsection*{NOTE GENERAL DATA}
{\small
\begin{longtable}{l>{\raggedright\arraybackslash}p{24mm}>{\raggedright\arraybackslash}p{58mm}>{\raggedright\arraybackslash}p{36mm}}
\toprule
no. & bits & parameter & values \\
\midrule\endfirsthead
\toprule no. & bits & parameter & values \\
\midrule\endhead
\bytes{198--1A4} & 6--0 & NAME (13 characters) & \bytes{20-7F} \\
\bytes{1A5} & 5, 3--0 & ATTRIBUTE FLAG: KEY OFF MODE, TONE ON/OFF & \bytes{00-01} --- two bits per tone, 2nd then 1st \\
\bytes{1A6} & 7, 6--0 & SAME GROUP DUMP: DUMP FLAG, GROUP NUMBER & \bytes{00-7F} --- flag 1 dumps the same group's sound \\
\bytes{1A7} & 3--0, 7--4 & MIXER OUTPUT SELECT: MAIN, SUB & \bytes{00-05} --- 0 no output, 1 MAIN, 2-4 SUB1-3, 5 EFFECT2 \\
\bytes{1A8} & 6--0 & EFFECT SEND VOLUME: EFFECT1 & \bytes{00-7F} \\
\bytes{1A9} & 6--0 & EFFECT SEND VOLUME: REVERB & \bytes{00-7F} \\
\bottomrule
\end{longtable}
}

\subsection*{NOTE TONE DATA}
{\small
\begin{longtable}{l>{\raggedright\arraybackslash}p{24mm}>{\raggedright\arraybackslash}p{58mm}>{\raggedright\arraybackslash}p{36mm}}
\toprule
no. & bits & parameter & values \\
\midrule\endfirsthead
\toprule no. & bits & parameter & values \\
\midrule\endhead
\bytes{1AA} & 5--0 & ATTRIBUTE DELAY & \bytes{00-32} --- 0 delay off \\
\bytes{1AB} & 7--0 & ATTRIBUTE PANNING & \bytes{00-80} --- 0-127 fixed pan, 128 random pan \\
\bytes{1AC} & 6--0 & TONE SELECT: TTN SELECT & \bytes{00-7F} \\
\bytes{1AD} & 7--6, 5--4, 3--0 & TONE KIND, TONE LOCATION, BANK SELECT & kind 0 NORMAL 1 DRUM 2 ATTACK \\
\bytes{1AE} & 7--0 & PITCH KEY SHIFT & \bytes{E8-18} --- -24 to +24 \\
\bytes{1AF} & 7--0 & PITCH DETUNE & \bytes{80-7F} --- -128 to +127 \\
\bytes{1B0} & 6--0 & LEVEL VOLUME & \bytes{00-7F} --- one step is 0.375 dB \\
\bytes{1B1} & 7--0 & LEVEL TOUCH DEPTH & \bytes{CE-32} --- -50 to +50 \\
\bytes{1B2} & 7--5 & LEVEL TOUCH CURVE & \bytes{00-06} --- 0 is -3 curve through 6 is +3 curve \\
\bytes{1B3--1B8} & 6--0 & LEVEL ENVELOPE SHAPE: attack time, decay1 time, sustain1 level, decay2 time, sustain2 level, release time & \bytes{00-64} --- one byte each, in that order \\
\bytes{1B9} & 7--0 & LEVEL ENV TOUCH FOLLOW ATTACK & \bytes{CE-32} \\
\bytes{1BA} & 7--0 & LEVEL ENV TOUCH FOLLOW DECAY & \bytes{CE-32} \\
\bytes{1BB} & 2--0, 7--5 & FILTER MODE, TOUCH CURVE & \bytes{00-06} --- 0 THROUGH, 1 LPF12+EQ, 2 HPF12+EQ, 3 LPF24, 4 HPF24, 5 BPF \\
\bytes{1BC} & 5--0 & FILTER TOUCH DEPTH & \bytes{00-32} \\
\bytes{1BD--1C0} &  & FILTER PARAMETER VALUE1-VALUE4 & meaning depends on the filter mode \\
\bottomrule
\end{longtable}
}

\subsection*{NOTE MODELING DATA}
{\small
\begin{longtable}{l>{\raggedright\arraybackslash}p{24mm}>{\raggedright\arraybackslash}p{58mm}>{\raggedright\arraybackslash}p{36mm}}
\toprule
no. & bits & parameter & values \\
\midrule\endfirsthead
\toprule no. & bits & parameter & values \\
\midrule\endhead
\bytes{1D8--1DA} & 6--0 & DRIVER THRESHOLD POINT: between A and B, B and C, C and D & \bytes{00-7F} --- provided that A < B < C < D \\
\bytes{1DB} & 6--0 & DRIVER A TTN SELECT & \bytes{00-7F} \\
\bytes{1DC} & 7--6, 5--4, 3--0 & DRIVER A KIND, LOCATION, BANK SELECT & kind 0 NORMAL 1 DRUM 2 ATTACK \\
\bytes{1DD--1DE} &  & DRIVER B: TTN select; kind, location, bank select &  \\
\bytes{1DF--1E0} &  & DRIVER C: TTN select; kind, location, bank select &  \\
\bytes{1E1--1E2} &  & DRIVER D: TTN select; kind, location, bank select &  \\
\bytes{1E3} & 7--6, 5--0 & RESONATOR GENERAL: INTERACTION CONNECT, RESONATOR TYPE & \bytes{00-3F} --- type 0 ORIGINAL, 1-63 preset \\
\bytes{1E4} &  & (RESERVE) &  \\
\bytes{1E5} & 7--0 & POSITION PARAMETER POSITION & \bytes{00-FA} --- 0.2 per cent per step, 0 to 50 per cent \\
\bytes{1E6} & 7, 6--0 & FORMANT FLAG, DEPTH & \bytes{00-7F} --- 0 FIXED, 1 MOVE \\
\bytes{1E7} &  & (RESERVE) &  \\
\bytes{1E8} & 7--0 & POSITION TOUCH & \bytes{CE-32} --- -50 to +50 \\
\bytes{1E9} & 5--0 & POSITION MOVEMENT WIDTH & \bytes{00-32} \\
\bytes{1EA} & 5--0, 7 & POSITION MOVEMENT SPEED, SAMPLE/HOLD & \bytes{00-32} \\
\bytes{1EB} & 6--0 & INTERACTION GAIN & \bytes{00-7F} \\
\bytes{1EC} &  & (RESERVE) &  \\
\bytes{1ED} & 7, 6--0 & MAIN RESONATOR MODE, FITTING & \bytes{00-7F} --- 0 positive feedback, 1 negative feedback \\
\bytes{1EE} & 7, 6--0 & MAIN RESONATOR SCALE, MUTING & \bytes{00-7F} \\
\bytes{1EF} & 7--0 & MAIN TOUCH FOLLOW TOUCH FITTING & \bytes{CE-32} \\
\bytes{1F0} & 7--0 & MAIN TOUCH FOLLOW TOUCH MUTING & \bytes{CE-32} \\
\bytes{1F1--1F3} & 6--0 & MAIN KEY FOLLOW: center key, low key limit, high key limit & \bytes{00-7F} \\
\bytes{1F4} & 7--0 & MAIN KEY FOLLOW SLOPE & \bytes{CE-32} \\
\bytes{1F5} & 7--0 & MAIN PITCH KEY SHIFT & \bytes{C4-3C} --- -60 to +60 \\
\bytes{1F6} & 7--0 & MAIN PITCH DETUNE & \bytes{80-7F} --- -128 to +127 \\
\bytes{1F7} & 7, 6--0 & SUB RESONATOR MODE, FITTING & \bytes{00-7F} \\
\bytes{1F8} & 7, 6--0 & SUB RESONATOR SCALE, MUTING & \bytes{00-7F} \\
\bytes{1F9} & 7--0 & SUB GAIN & \bytes{9C-64} \\
\bytes{1FA} & 7--0 & SUB TOUCH FOLLOW TOUCH FITTING & \bytes{CE-32} \\
\bytes{1FB} & 7--0 & SUB TOUCH FOLLOW TOUCH MUTING & \bytes{CE-32} \\
\bytes{1FC} & 7--0 & SUB TOUCH FOLLOW TOUCH SUB GAIN & \bytes{CE-32} \\
\bytes{1FD--1FF} & 6--0 & SUB KEY FOLLOW: center key, low key limit, high key limit & \bytes{00-7F} \\
\bytes{200} & 7--0 & SUB KEY FOLLOW SLOPE & \bytes{CE-32} \\
\bytes{201} & 7--0 & SUB PITCH KEY SHIFT & \bytes{C4-3C} \\
\bytes{202} & 7--0 & SUB PITCH DETUNE & \bytes{80-7F} \\
\bottomrule
\end{longtable}
}



\section{Drum sound}
\textsc{drum sound} is reached at \bytes{18 00 00}, and messages addressed there are
accepted. Its contents are not enumerated in this edition.
